\documentclass[12pt]{article}
\usepackage[T1]{fontenc}
\usepackage[english]{babel} % english as language
\usepackage{csquotes}
\usepackage[onehalfspacing]{setspace} % 1.5 Zeilenabstand
\usepackage{parskip} % for paragraph line skip
\usepackage{fancyhdr} % for headers
\fancyhf{}
\fancyhead[C]{\nouppercase{\firstmark}}
\renewcommand{\headrulewidth}{0pt} % change line under header
\renewcommand{\footrulewidth}{0pt}
\pagestyle{fancy}
\setlength{\headheight}{15.04742pt}
\fancypagestyle{maintext}{ % pagestyle for main text for page numbers
    \fancyhf{}
    \fancyhead[C]{\nouppercase{\firstmark}}
    \fancyfoot[R]{\thepage}
    \renewcommand{\headrulewidth}{0pt}
    \renewcommand{\footrulewidth}{0pt}
}
% Update \sectionmark to show first section in header
\renewcommand{\sectionmark}[1]{%
    \markboth{\thesection\ #1}{}
}
\renewcommand{\subsectionmark}[1]{}

% for tables/figures
\usepackage{graphicx}
\usepackage{caption}
\usepackage{array}
\graphicspath{ {figures/} } % set path to figures
\renewcommand{\listfigurename}{Abbildungsverzeichnis}
\renewcommand{\listtablename}{Tabellenverzeichnis}

% Appendix
\usepackage[toc,page]{appendix}
\renewcommand\appendixname{Appendix}
\renewcommand\appendixpagename{Appendix}
\renewcommand\appendixtocname{Appendix}


% \usepackage{sectsty} -> falls section font size 14 haben muss
% \sectionfont{\large\bfseries}

% use geometry for page layout
\usepackage[a4paper, left=2.5cm, right=2.5cm, top=2.5cm, bottom=2cm]{geometry}

% Citation: BibLaTeX als author-year style, benötigt xelatex->biber->xelatex(2x) als recipe!
\usepackage[
backend=biber,
style=authoryear]{biblatex}
\addbibresource{refs.bib}

% Start des Dokuments
\begin{document}

% Add centered Title page
\begin{titlepage}
    \begin{center}
        \vspace*{0.5cm}
            
        \LARGE
        \textbf{Do university instructors spot AI?}
            
        \vspace{0.2cm}
        \Large
        Investigating the Detectability of AI-Generated\\
        Student Writing in Higher Education
            
        \vspace{2cm}
            
        \textbf{Lukas Schaffrath}\\
        Matriculation-Nr. 411793\\
        M.A. Cognitive, Digital and\\
        Empirical English Studies
            
        \vspace{2cm}
            
        Seminar: Empirical Research - Digital Humanities\\
        Instructor: Dr. Elma Kerz
            
        \vspace{2cm}
            
            
        \Large
        English and American Studies\\
        RWTH Aachen University\\
        Winter term 2025\\
        30. March 2025
            
    \end{center}
\end{titlepage}

% Add table of contents with emtpy page number
\addtocontents{toc}{\protect\thispagestyle{empty}}
\tableofcontents

% List of figures and tables
\newpage
\thispagestyle{empty}
\listoffigures
\vspace{0.5cm}
\listoftables

\newpage
\thispagestyle{empty}
\begin{abstract}
The rapid diffusion of generative artificial intelligence (GenAI) in higher education has raised substantial concerns regarding academic integrity and the detectability of AI-generated student writing. While prior research suggests that educators often fail to distinguish synthetic texts from human-authored ones, little is known about the detection capabilities of subject-matter experts. This exploratory study investigates whether university instructors in English Studies can reliably differentiate between human-written and AI-generated academic conclusions.
Using an online binary classification task, four experienced university instructors evaluated ten standardized conclusion excerpts, half written by students and half generated by a large language model. Results indicate a mean classification accuracy of 80\%, significantly exceeding chance level. However, performance was asymmetrical: instructors demonstrated near-perfect recognition of authentic student texts, while AI-generated samples proved considerably more difficult to detect.
These findings suggest that expert instructors possess a refined sensitivity to authentic student writing, yet remain vulnerable to increasingly “humanized” AI outputs. The study highlights the ethical importance of high specificity in AI detection and underscores the need for enhanced AI literacy, process-oriented assessment designs, and a pedagogical shift toward critical engagement with generative technologies in higher education.
\end{abstract}

% New page, page count to 1
\newpage
\setcounter{page}{1}
\pagestyle{maintext}

% Add Main Text
\section{Introduction}
Since the emergence of widely accessible Artificial Intelligence (AI) text generation tools, higher education has undergone a transformation affecting teaching, learning, as well as assessment practices. The rapid advancement and integration of large language models (LLMs), such as OpenAI's ChatGPT, have fundamentally altered the educational landscape, challenging long-established assumptions about authorship, originality, and academic performance \autocite[2]{fleckenstein_teachers_2024}. Within the broader context of the digital transformation of education, AI is increasingly described as a revolutionary technological force, as it enables machines to perform tasks that typically require human intelligence, sometimes considered exclusive to human intellectual activity \autocite[92]{kumar_ai_2025}.

In educucational contexts, this technological development can be associated with considerable pedagogical potential, as Scholars highlight opportunities for personalized learning environments, adaptive feedback mechanisms, and enhanced educational quality through instruction tailored to individual student needs \autocites[449]{orhani_risks_2025}[96]{kumar_ai_2025}. From this perspective, AI tools could be seen as a supportive tool enhancing inclusivity and supporting diverse student populations by lowering cognitive barriers to academic participation. 

However, the widespread availability of LLMs also raises significant concerns regarding academic integrity, as students may exploit these tools to generate written work, potentially undermining the authenticity of their submissions \autocite[2]{fleckenstein_teachers_2024}. Students relying on AI-generated texts may gain an unfair advantage in grading contexts, while bypassing essential learning processes, e.g. critical thinking and academic writing ompetence \autocites[2]{fleckenstein_teachers_2024}[4]{rahimi_ai_2025}. As a result, maintaining and educating validity and fairness of academic assessments, as well as scientific integrity, have become a primary challenge for higher education institutions. Therefore, instructors must be able to distinguish between original student work and automated outputs, generated by AI, to ensure fair evaluation \autocites[2]{fleckenstein_teachers_2024}[3]{deriba_fairness_2025}. Central to this challenge is the question of detectability. Previous research points to significant gaps in instructors' ability to recognize AI-generated texts, sometimes performing no better than chance level, with accuracy rates around 50\% \autocite[1]{fleckenstein_teachers_2024}. Particularly concerning is the documented overconfidence of instructors in their detection abilities, despite demonstrating high error rates \autocite[1]{fleckenstein_teachers_2024}. While technological detection approaches based on machine-learning models, such as BERT, can achieve classification accuracies of beyond 90\% under experimental conditions, these systems remain largely inaccessible to instructors in everyday grading routines \autocite[6]{walker_text_2024}. Furthermore, this problem is worsened by the increased sophistication of AI-generated texts, which can closely mimic human writing by usage of humanizing prompts to imitate student-level linguistic features. 

Given these challenges, there is an urgent need for empirical research examining the detection capabilities of experts in higher education, where written texts not only serve as a medium for knowledge demonstration and assessment, but also play a crucial role in the development of academic practices. Addressing this gap, the present study tries to investigate whether experienced university instructors are able to distinguish between student-written and AI-generated texts. The guiding research question of this study is therefore: \textit{To what extent can university instructors distinguish between AI-generated and human-written student texts?}. Based on this, two competing hypotheses can be formulated. The null hypothesis (H0) assumes that instructors' detection accuracy does not differ significantly from the 50\% chance level. In contrast, the alternative hypothesis (H1) assumes that instructors are able to identify AI-generated texts with an accuracy significantly higher than chance. By empirically testing these hypotheses, this study aims to contribute to the broader discourse on academic integrity, assessment validity, and the practical implications of GenAI for higher education. 

\newpage
\section{Context and Background}

The rapid development in usage of generative artificial intelligence in higher education has exposed fundamental limitations in human text evaluation, even among trained professionals. Empirical research indicates that humans struggle to distinguish between AI-generated and human-written content, particularly in educational settings where texts follow standardized academic conventions. Fleckenstein et al.\ demonstrated that both novice teachers ($N = 89$) and experienced teachers ($N = 200$) were largely unable to identify AI-generated texts from authentic student essays, with performance frequently approximating chance level \autocite[1]{fleckenstein_teachers_2024}. Moreover, pre-service teachers correctly identified only 45.1\% of AI-generated texts, while experienced teachers correctly identified merely 37.8\% of synthetic samples, indicating that professional experience alone does not guarantee improved detection accuracy (ebd.).

A major contributing factor to this limitation is the so-called \emph{overconfidence effect}. Instructors tend to overestimate their diagnostic abilities, particularly when they assume a text to be student-written, which leads to misclassification despite high subjective certainty \autocite[1]{fleckenstein_teachers_2024}. This discrepancy between perceived and actual competence highlights a critical expertise gap in education assessment practices, as evaluators remain unaware of the extent to which AI-generated texts evade human judgment.

This problem is further worsened by the deceptive nature of modern LLMs. Contemporary AI can be prompted to produce humanized outputs resembling student writing by incorporating typical learner deficiencies, such as simplified syntax, vague argumentation, hedging strategies, or even minor grammatical and spelling errors. As a result, experienced instructors often fail to recognize that AI systems are capable of intentionally generating weak or mediocre texts. Instead, such outputs are frequently misattributed to human students. Fleckenstein et al.\ report that AI-generated texts are often assessed as being of higher quality than genuine student work, even when evaluators correctly identify the source, suggesting that perceived quality and source attribution are not consistently aligned \autocite[2]{fleckenstein_teachers_2024}. This stealth capability enables students to deceive instructors while maintaining plausible deniability.

In response to these challenges, many instructors may have turned to automated AI-detection software, such as GPTZero, promising high-accuracy AI content-detection from many LLMs \autocite{noauthor_gptzero_nodate}. However, research suggests that these tools are currently not reliable enough to serve as definitive instruments for identifying AI-generated content. Chaka shows that existing detection systems are \enquote{not yet fully ready} to accurately and consistently detect AI-generated texts across different linguistic and contextual conditions \autocite[10]{chaka_detecting_2023}. Detection efficacy is further undermined by simple manipulation techniques, including paraphrasing, the insertion of extraneous characters, or the use of machine translation tools, all of which can cause AI-generated texts to be misidentified as human-produced \autocite[8--9]{chaka_detecting_2023}. Moreover, automated detectors typically rely on probabilistic estimations rather than verifiable evidence, rendering them unsuitable as conclusive proof in cases of suspected academic misconduct.

In contrast to both human evaluators and commercially available detection software, specialized technical models demonstrate substantially higher robustness. BERT-based classification systems, utilizing contextual embeddings to capture subtle semantic and linguistic patterns, have been shown to achieve accuracy rates exceeding 92\% in distinguishing AI-generated from human-written texts \autocite[6]{walker_text_2024}. By modeling deep contextual relationships remaining largely invisible to human readers, these systems establish an important technical benchmark for detection performance. Despite their high accuracy, though, such computational models remain largely inaccessible to university instructors in everyday grading and assessment practices, limiting their practical applicability in educational contexts.

Beyond technical and evaluative challenges, the integration of AI in higher education raises broader socio-pedagogical concerns. Several studies emphasize a critical lack of AI literacy among educators, who often lack the formal training necessary to understand the functioning, limitations, and ethical implications of AI-based systems \autocite[3]{deriba_fairness_2025}. This lack of preparation not only complicates the detection of synthetic content but also increases the risk of ethical misjudgments, including false accusations of academic misconduct. Furthermore, scholars warn that an uncritical reliance on AI may contribute to the mechanization of education, potentially diminishing students' creativity, critical thinking, and intellectual autonomy \autocite[4]{rahimi_ai_2025}.

To address these challenges, several authors advocate for a shift from a predominantly coding-oriented approach to a decoding-oriented approach in AI education. Rather than focusing solely on the technical use or development of AI systems, this perspective emphasizes the importance of teaching both educators and students how to critically analyze, interpret, and deconstruct AI-generated outputs \autocite[5]{rahimi_ai_2025}. From an ethical standpoint, maintaining high specificity, meaning the ability to correctly identify authentic human work, is of particular importance, as false positives may unjustly penalize students and undermine trust in academic assessment systems \autocite[449]{orhani_risks_2025}. Consequently, understanding the limitations of both human and automated detection mechanisms is essential for developing fair, transparent, and pedagogically sound strategies for assessment in the age of generative AI.

\newpage
\section{Methodology}

This study employed an exploratory online experimental design with a binary classification task to investigate the detectability of synthetic academic texts. The primary objective was to measure the ability of university instructors to distinguish between authentic student-written conclusions and conclusions generated by a LLM under controlled conditions. The experimental framework was designed to capture not only decision accuracy but also confidence and response time as complementary indicators of evaluative behavior (see Anhang!!).

\subsection{Participants}

The participant sample consisted of $N = NUMBER OF PARTICIPANTS$ university instructors recruited from specialized academic contexts, including English Linguistics, Language Competence, and the Department of English, American and Romance Studies (IfAAR) at RWTH Aachen University (see Appendix). All participants possessed experience in assessing academic writing, with a mean teaching experience of MEAN YEARS ($SD = STANDARD DEVIATION$), ranging from 2 to 10 years. The inclusion of instructors from linguistics- and language-focused disciplines was intended to test the upper limits of human detection performance, as these experts routinely evaluate complex written texts at an advanced academic level.

\subsection{Materials}

The data comprised a total of ten short academic text excerpts, divided evenly into a human-written corpus and an AI-generated corpus. The human-written texts consisted of five excerpts drawn from authentic student term papers submitted at RWTH Aachen University between 2021 and 2024 (see Human Texts, p.). The AI-generated texts consisted of five corresponding excerpts generated using the GPT~5.1 model (see AI Generation Log).

To ensure structural comparability, all ten excerpts were standardized to a length of exactly eight sentences and were taken exclusively from the beginning of the \emph{Conclusion} section of the academic papers. This section was chosen because it typically integrates argumentation, evaluation, and summarization, making it particularly suitable for assessing authorial style and coherence. To maximize task difficulty and validity, the AI model was prompted to produce texts that resembled realistic student-level academic writing. The prompting strategy explicitly instructed the model to adopt a natural academic tone characteristic of B.A.\ and M.A.\ level submissions, including the use of hedging expressions, tentative phrasing, and minor stylistic inconsistencies commonly observed in student writing (see Prompt Guidelines, p.~1; see AI Output Examples, p.~1).

\subsection{Procedure}

Data collection and exportation was conducted via a custom-built web-based interface hosted on \textit{https://lstoneyy.pythonanywhere.com}, which ensured a standardized and controlled presentation of stimuli across participants. Each participant was presented with all ten text excerpts in a fully randomized order to control for potential order, or learning effects.
For each excerpt, participants were required to complete two tasks. First, they performed a binary classification by indicating whether they believed the text to be human-written or AI-generated. Second, they provided a confidence rating for each decision on a five-point Likert scale ranging from 1 (\emph{very unsure}) to 5 (\emph{very sure}), allowing for the assessment of calibration between confidence and accuracy (see Procedure Overview, p.~1).

In addition to these self-reported measures, the system automatically recorded response times in milliseconds for each classification decision. Response time was used as an index of cognitive effort and deliberation during the evaluative process (see Statistics Appendix, p.~2; see Results, p.~9).

\subsection{Analytical Framework}

Upon data input on the used website, each collected entry was automatically saved in a reponse data-table, comprised of the participant ID, the presented text ID, the binary classification decision, the confidence rating, and the recorded response time in milliseconds (see Data Sheet Example, p.~1). By using a framework to export the response data as a structured .csv-file, the entire analysis pipeline from data collection to statistical evaluation was streamlined and standardized.

All statistical analyses were then performed using a Python script, which is provided in the Appendix. Descriptive statistics, inferential tests, and visualizations were generated using standard libraries, including NumPy, SciPy, and Matplotlib. It not only conducted the statistical analysis, but exported plots, text ID and participant ID references as text files, and all results in a markdown-file (see OUTPUT). The automated script was developed for reproducibility and transparency, allowing for independent verification of results. It can be used with any .csv-table, though modifications may be necessary depending on the structure of the input data (see Script, p.~1). 

The primary dependent variable was classification accuracy, operationalized as the proportion of correctly identified texts across all trials and participants. Accuracy was analyzed both at the aggregate level and at the individual participant level, as can be seen in chapter 4. To test the central research hypothesis, a one-sample $t$-test was conducted to compare the observed mean accuracy against the theoretical chance level of 50\%. 
To further decompose classification performance, a confusion matrix was constructed to calculate sensitivity (true positive rate for AI-generated texts) and specificity (true negative rate for correctly identifying human-written texts). These measures were included to assess not only overall accuracy but also the ethical dimension of misclassification, particularly the risk of falsely identifying authentic student work as AI-generated.

Finally, Pearson correlation coefficients ($r$) were computed to examine the relationships between classification accuracy, confidence ratings, response times, and years of teaching experience. This analysis aimed to explore potential associations between expertise, metacognitive calibration, and detection performance.

\newpage
\section{Discussion and Results}

\subsection{Overall Detection Performance and Hypothesis Testing}

The statistical analysis revealed a mean classification accuracy of 80.00\% ($SD = 16.33\%$) across all participating university instructors (see Results, p.~12). To evaluate whether this performance exceeded chance level, a one-sample $t$-test was conducted against the theoretical baseline of 50\%. The analysis yielded a statistically significant result, $t(3) = 3.674, p = 0.0349$, leading to the rejection of the null hypothesis ($H_0$) that instructors are unable to distinguish between human-written and AI-generated texts beyond random guessing (see Results, p.~13).

These findings indicate that subject-matter experts in higher education possess a systematic ability to differentiate between text sources. This result stands in marked contrast to previous studies in which both novice and experienced teachers performed at or near chance level when evaluating AI-generated student essays \autocite[1]{fleckenstein_teachers_2024}. The magnitude of this effect is further underscored by a Cohen’s $d$ of 1.837, which constitutes a large effect size and suggests robust discrimination capabilities within this specialized sample (see Results, p.~13).

\subsection{The Asymmetry of Recognition: Specificity versus Sensitivity}

A more fine-grained analysis revealed a pronounced asymmetry in detection performance depending on text origin (see Results, p.~12). Instructors demonstrated near-perfect recognition of authentic student writing, correctly identifying 19 out of 20 human-written excerpts, corresponding to an accuracy rate of 95\%. In contrast, AI-generated texts proved substantially more difficult to detect: only 65\% of synthetic excerpts were correctly classified, indicating that more than one third of AI-generated content went undetected.

This imbalance between specificity and sensitivity has important ethical implications. The high specificity rate of 95\% significantly reduces the risk of false-positive classifications, that is, the erroneous attribution of AI authorship to genuine student work. From a pedagogical and fairness-oriented perspective, this finding is particularly salient, as false accusations of AI misuse may disproportionately harm students and undermine trust in academic assessment practices \autocite[11]{deriba_fairness_2025}. Consequently, while imperfect AI detection remains a concern, the instructors’ strong ability to correctly recognize human writing constitutes a critical safeguard for academic integrity.

\newpage
\subsection{Metacognitive Calibration and Instructor Confidence}

Beyond accuracy, the study examined instructors’ metacognitive awareness as reflected in their confidence ratings. Participants reported a mean confidence level of 3.65 on a five-point Likert scale (see Results, p.~11). Importantly, confidence varied systematically with decision accuracy. Instructors expressed significantly higher confidence for correct classifications ($M = 3.81$) than for incorrect ones ($M = 3.00$), a difference that reached statistical significance ($p = 0.0038$; see Results, p.~13).

Moreover, confidence and accuracy were strongly positively correlated ($r = 0.707$), suggesting a high degree of metacognitive calibration. This finding contrasts sharply with earlier research documenting an overconfidence effect, whereby instructors substantially overestimated their ability to detect AI-generated texts despite poor objective performance \autocite[8]{fleckenstein_teachers_2024}. In the present study, participants appeared to possess a more realistic understanding of their own diagnostic limits, which may be attributable to their subject-matter expertise and frequent engagement with advanced academic writing.

\subsection{Item-Level Analysis and the GPT~5.1 ``Stealth'' Effect}

An item-level analysis revealed considerable variation in detectability across individual text excerpts. While most human-written texts were identified with perfect accuracy, one AI-generated excerpt (Text~6) emerged as particularly deceptive, achieving a classification accuracy of only 25\% (see Results, p.~14). This text was generated using a prompt that emphasized reflective content, including personal emotions and subjective experiences related to ``digital diaries.''

The low detectability of this item suggests that LLMs may be especially effective at evading expert judgment when operating within highly subjective or introspective genres. Such genres appear to align closely with instructors’ expectations of authentic student voice, thereby blurring the boundary between human and synthetic authorship. This finding supports earlier concerns that so-called ``humanizing'' prompts can successfully mimic linguistic features that evaluators associate with genuine student writing \autocite[10]{chaka_detecting_2023}. As LLMs continue to improve in simulating stylistic nuance, these stealth effects may become increasingly pronounced.

\subsection{Cognitive Processing and the Role of Intuition}

Response time analysis provided further insight into the cognitive processes underlying detection performance. A strong negative correlation was observed between response time and accuracy ($r = -0.933, p = 0.0670$; see Results, p.~13), indicating that faster judgments were often more accurate than slower, more deliberative ones. Although this correlation narrowly missed conventional thresholds for statistical significance, the strength of the association suggests a meaningful trend.

This pattern implies that expert instructors may rely on rapid, intuitive assessments when evaluating text authenticity, whereas extended analytical processing may introduce doubt or lead to the overinterpretation of ambiguous stylistic cues. Such findings align with dual-process theories of cognition, which posit that expert performance often depends on fast, experience-based heuristics rather than exhaustive analytical reasoning.

\subsection{Contextualizing the Findings: Expertise versus Technology}

When situated within the broader literature, the instructors’ mean accuracy of 80\% clearly exceeds the generalist baseline reported in prior studies but remains below the 90\% and higher accuracy rates achieved by specialized technical detection models such as BERT-based classifiers \autocite[6]{walker_text_2024}. However, unlike human evaluators, automated detection tools are highly vulnerable to simple manipulation strategies, including paraphrasing or machine translation, which can substantially reduce detection reliability \autocite[2]{chaka_detecting_2023}.

In this context, the instructors’ high specificity emerges as a crucial advantage. While neither human nor automated detection offers a perfect solution, expert judgment remains the most reliable mechanism for protecting students from unjust accusations. These findings reinforce calls for a shift toward a ``decoding'' approach in higher education, in which instructors leverage their disciplinary expertise not merely to detect AI use, but to guide students in critically understanding and reflecting on academic writing processes in the age of generative AI \autocites[4]{rahimi_ai_2025}[96]{kumar_ai_2025}.

\newpage
\section{Conclusion}

This study set out to investigate whether university instructors are able to distinguish between human-written and AI-generated academic texts beyond chance level. The empirical findings demonstrate that the participating instructors achieved an overall classification accuracy of 80.00\% (see Results, p.~15). Statistical testing confirmed that this performance was significantly higher than the theoretical chance level of 50\% ($p = 0.0349$), providing sufficient evidence to reject the null hypothesis that instructors are unable to differentiate between text origins beyond random guessing (see Results, p.~16).

At the same time, the analysis revealed a marked asymmetry in detection performance. While instructors demonstrated near-perfect recognition of authentic student writing, achieving a specificity rate of 95\%, their sensitivity toward AI-generated texts was substantially lower at 65\% (see Results, p.~15). This imbalance suggests that although instructors are highly reliable in avoiding false accusations of AI misuse, a considerable proportion of synthetic texts may still pass as human-authored.

In contrast to earlier research reporting chance-level performance among educators \autocite[1]{fleckenstein_teachers_2024}, the present findings indicate that subject-matter experts in English Studies possess robust discrimination skills. The high specificity observed in this study suggests that experienced instructors have developed a refined sensitivity to features associated with an authentic student voice. From an ethical standpoint, this constitutes a crucial safeguard against false-positive classifications, which may otherwise undermine academic fairness and disproportionately disadvantage students \autocite[11]{deriba_fairness_2025}. However, the particularly low detectability of one AI-generated excerpt—correctly classified in only 25\% of cases—demonstrates that modern language models can effectively imitate reflective and emotionally inflected student writing styles when guided by carefully designed prompts (see Results, p.~17; \autocite[10]{chaka_detecting_2023}).

These findings carry important pedagogical and ethical implications for higher education. First, they underline the limitations of relying exclusively on human intuition for AI detection. While expert judgment remains a valuable filter, instructors require formal AI literacy training to better understand the evolving capabilities and limitations of generative models \autocites[426]{orhani_risks_2025}[109]{kumar_ai_2025}. Second, the results support calls for a shift from a purely “coding” approach—focused on identifying AI-generated text—to a “decoding” approach, in which students and instructors collaboratively analyze and reflect on AI outputs as part of the learning process \autocite[188]{rahimi_ai_2025}. Such an approach may foster critical engagement with academic writing rather than incentivizing concealment strategies.

From an assessment perspective, the study reinforces the need to rethink traditional text-based assignments. To safeguard academic integrity, institutions may benefit from prioritizing low-stakes diagnostic writing tasks, process-oriented assessments, and oral examinations that test a student’s conceptual understanding rather than solely the final written product \autocites[51]{fleckenstein_teachers_2024}[521]{orhani_risks_2025}. These measures may help reduce the incentives for AI misuse while preserving the pedagogical value of academic writing.

Several limitations must be acknowledged. Most notably, the small sample size ($N = 4$) restricts the generalizability of the findings beyond the specific institutional and disciplinary context of this study (see Results, p.~14). Additionally, the experimental design relied on short, decontextualized text excerpts evaluated without reference to a student’s broader writing history, which may not fully reflect authentic grading situations (see Results, p.~18). Future research should therefore employ larger and more diverse participant samples and examine instructor performance against newer generations of generative AI models as they continue to evolve.

In conclusion, university instructors are not merely guessing when confronted with AI-generated academic writing, but they are increasingly challenged by highly sophisticated systems designed to mimic human authorship. While human expertise remains an indispensable component of academic assessment, the continued viability of the traditional text-based essay as a secure evaluation format appears uncertain unless it is complemented by reflective, process-oriented, and supervised forms of assessment \autocites[51]{fleckenstein_teachers_2024}[520]{orhani_risks_2025}.


% print bibliography with heading in table of contents
\newpage
\printbibliography[heading=bibintoc]
\pagestyle{empty}

% Starte Anhang
\newpage
\begin{appendices}

    % Sektionen
    % \section{Response Data Sheet Example}
    % \begin{center}
    %     % \includegraphics[height=14cm, angle=90, keepaspectratio]{static/Unterrichtsstunde_1_Tabelle.pdf}
    %     \captionof{table}{Tabellarischer Verlaufsplan zur ersten Unterrichtsstunde}
    % \end{center}
    %
    \section{Declaration of Academic Integrity}
    \begin{center}
        \includegraphics[height=22cm, angle=0, keepaspectratio]{static/declaration.pdf}
    \end{center}
\end{appendices}

\end{document}
